\section{Prefix-Aware Recertification}
\label{sec:theory}

Let a plan $P$ be certified and executing. The \emph{frontier} $\phi$
records, for each workflow $w$, the number $\phi(w)$ of its pieces that have
committed. A workflow $r$ proposes to replace its uncommitted suffix with a
new sequence of pieces, giving a revised plan $P'$. The runtime must decide
whether to admit $P'$ and, if the proposed chopping is too fine, whether it
can be repaired.

\subsection{Soundness of recertification}

\begin{theorem}[Inherited isolation]
\label{thm:inherit}
If $P'$ is SC-cycle-free, its committed prefix is identical to the prefix
actually executed, and the runtime satisfies E1--E4, then the resulting
history is conflict-serializable at workflow granularity.
\end{theorem}

\begin{proof}
The pieces that already ran are, by construction, exactly the prefix pieces of
$P'$ with the footprints they declared, and E3 guarantees that their actual
accesses were inside those footprints. The observed execution is therefore a
legal execution of the chopping $P'$ under the execution rules. The classical
chopping theorem~\cite{shasha1995chopping} applies to $P'$ and yields
equivalence to a serial execution of the workflows of $P'$.
\end{proof}

Theorem~\ref{thm:inherit} is an application of established theory to a revised
plan, not a new isolation theorem. Its force is in the hypothesis:
recertification must cover the \emph{whole} batch, including pieces that have
committed. Certifying only the new suffix discards exactly the conflicts that
close a cycle through the past.

\subsection{What repairs remain after commitment}

Before execution, a rejected chopping is repaired by merging pieces. After
commitment only the uncommitted suffix may be regrouped.

\begin{definition}
$P''$ is a \emph{future-only coarsening} of $P'$ at $\phi$ if the first
$\phi(w)$ pieces of every workflow $w$ are unchanged and the remaining pieces
of $w$ are grouped into contiguous runs. The \emph{sealing}
$\seal(P',\phi)$ is the future-only coarsening that merges each workflow's
entire remainder into a single piece.
\end{definition}

\begin{lemma}[Sealing is coarsest]
\label{lem:coarsest}
Every future-only coarsening of $P'$ at $\phi$ is a refinement of
$\seal(P',\phi)$.
\end{lemma}

\begin{proof}
Both agree on the committed prefix. Sealing places each workflow's whole
remainder in one piece; any other future-only coarsening partitions that same
remainder into contiguous runs, which is a refinement of the single piece.
\end{proof}

\begin{lemma}[Cycles persist under refinement]
\label{lem:persist}
If a chopping contains a mixed simple cycle, so does every refinement of it.
\end{lemma}

This is the classical non-removability property of
chopping~\cite{shasha1995chopping}: splitting a piece adds S~edges among the
parts and distributes the piece's C~edges among them, so a mixed cycle through
the original piece survives as a mixed cycle through its parts.

\begin{theorem}[Admissibility characterization]
\label{thm:char}
A revision is admissible under future-only coarsening if and only if
$\seal(P',\phi)$ is SC-cycle-free.
\end{theorem}

\begin{proof}
($\Leftarrow$) If the sealing is SC-cycle-free, it is itself a certified
future-only coarsening, so the revision is admissible.
($\Rightarrow$) Suppose the sealing contains a mixed simple cycle. By
Lemma~\ref{lem:coarsest} every future-only coarsening is a refinement of the
sealing, and by Lemma~\ref{lem:persist} every such refinement contains a mixed
simple cycle. Hence no future-only coarsening certifies.
\end{proof}

Theorem~\ref{thm:char} turns admissibility into one test on the coarsest
candidate. It also tells the runtime that when that test fails, searching the
exponentially many intermediate choppings is pointless.

\begin{corollary}[Impossibility of future-only repair]
\label{cor:imp}
If $\seal(P',\phi)$ contains a mixed simple cycle, then no future-only repair
exists, and the runtime must refuse the revision or compensate the committed
prefix. Every such cycle passes through a committed piece.
\end{corollary}

\begin{proof}
The first claim is the contrapositive of Theorem~\ref{thm:char}. For the
second, in the sealed plan each workflow's uncommitted remainder is a single
piece, so every S~edge has at least one committed endpoint; a mixed cycle
contains an S~edge.
\end{proof}

\begin{figure}[t]
\centering
\begin{tikzpicture}[
  piece/.style={draw=ccmuted, rounded corners=2pt, minimum height=5mm,
                minimum width=13mm, font=\scriptsize, inner sep=1.5pt},
  done/.style={piece, fill=ccgrid},
  prop/.style={piece, densely dotted, line width=0.8pt},
  S/.style={ccblue, line width=0.9pt},
  C/.style={ccorange, line width=0.9pt, dashed},
  lab/.style={font=\scriptsize}]
% timeline
\node[lab, anchor=east] at (-0.2,1.2) {$W_0$};
\node[lab, anchor=east] at (-0.2,0.0) {$W_1$};
\node[done] (p0) at (0.7,1.2) {$r(x)$};
\node[done] (q)  at (2.6,0.0) {$w(x),\,r(y)$};
\node[prop] (p1) at (4.9,1.2) {$w(y)$};
\draw[ccmuted, line width=0.6pt] (3.75,1.7) -- (3.75,-0.5);
\node[lab, ccmuted, anchor=south] at (3.75,1.7) {frontier $\phi$};
\node[lab, ccmuted] at (1.6,-0.55) {committed};
\node[lab, ccmuted] at (4.9,-0.55) {proposed};
% edges
\draw[S] (p0) -- (p1) node[midway, above, lab, ccblue] {S};
\draw[C] (p0) -- (q)  node[midway, left=1pt, lab, ccorange] {C on $x$};
\draw[C] (q)  -- (p1) node[midway, right=2pt, lab, ccorange] {C on $y$};
\end{tikzpicture}
\caption{The impossibility witness of Corollary~\ref{cor:imp}. $W_0$ has
committed a read of $x$; $W_1$ has committed a piece that wrote $x$ and read
$y$. $W_0$ now proposes to write $y$. The cycle is mixed and passes through
$W_0$'s committed read. $W_0$'s remainder is already a single piece, so there
is nothing left to merge, and $W_1$ having finished does not remove its
edges.}
\label{fig:witness}
\end{figure}

\paragraph{Waiting does not help.}
The practically important part of Corollary~\ref{cor:imp} is that the refusal
cannot be waited out. One might hope that letting all other workflows finish
would dissolve the conflict. It does not. A workflow that has fully committed
still contributes its S~clique and its C~edges, because those interleavings
really happened. Figure~\ref{fig:witness} shows the smallest instance: $W_1$
is finished, $W_0$'s remainder is a single piece, and the cycle is still
there. The schedule it stands for is a real anomaly---$W_0$ read $x$ before
$W_1$ wrote it, so $W_0$ precedes $W_1$; $W_1$ read $y$ before $W_0$ writes
it, so $W_1$ precedes $W_0$. The runtime's options are to refuse the revision
or to run a compensating action for the committed prefix, which is an
application-level obligation rather than a graph property.

\subsection{Admission must be serialized}

\begin{theorem}[Serialized admission]
\label{thm:serial}
Two revisions that are each admissible against a common plan snapshot may be
jointly inadmissible. Admission decisions must therefore be evaluated and
installed one at a time against a single plan snapshot.
\end{theorem}

\begin{proof}
By counterexample. $W_0$ has committed a read of $x$ and $W_1$ a read of $y$.
$W_0$ proposes to write $y$ and $W_1$ proposes to write $x$. Against the
common snapshot, each revision alone adds one C~edge and creates no cycle, so
each is admitted. Applied together, the plan contains the cycle
$r_0(x)$--$w_0(y)$--$r_1(y)$--$w_1(x)$--$r_0(x)$, which alternates S and
C~edges, and the plan admits a non-serializable execution.
\end{proof}

The runtime therefore evaluates and installs revisions under a single
admission lock, which also fixes the frontier used for sealing. A piece that
is currently executing is treated as immutable, exactly like a committed one,
so that a concurrent admission cannot regroup a piece another thread is
running.

\subsection{The admission algorithm}

\begin{algorithm}[t]
\caption{Admitting a plan revision}
\label{alg:admit}
\begin{algorithmic}[1]
\Require current plan $P$, frontier $\phi$, proposed revision of workflow $r$
\Ensure an installed plan, or a refusal with a witness cycle
\State \textbf{acquire} the admission lock \Comment{Theorem~\ref{thm:serial}}
\State $P' \gets P$ with $r$'s uncommitted suffix replaced
\If{$P'$ is SC-cycle-free}
  \State \Return \textsc{Admit}$(P')$ \Comment{finest plan, unchanged}
\EndIf
\State $S \gets \seal(P',\phi)$
\If{$S$ has a mixed simple cycle $c$}
  \State \Return \textsc{Refuse}$(c)$ \Comment{Corollary~\ref{cor:imp}}
\EndIf
\Repeat \Comment{guided coarsening, suffix only}
  \ForAll{mixed biconnected blocks $b$ of $P'$}
    \State pick a workflow $w$ with an S~edge in $b$ beyond $\phi(w)$
    \State merge the interval of $w$ whose merge adds the least
      footprint~$\times$~duration
  \EndFor
\Until{$P'$ is SC-cycle-free}
\State \Return \textsc{AdmitCoarsened}$(P')$
\end{algorithmic}
\end{algorithm}

Algorithm~\ref{alg:admit} proceeds in three steps. It first certifies the
proposed plan as given, because when the proposal already certifies we want
the finest plan, not the coarsest. If that fails it tests the sealing; by
Theorem~\ref{thm:char} a failure here is final, and the algorithm returns the
witness cycle and the workflows it implicates so that the application can
choose compensation or an alternative plan. Otherwise it runs a guided
coarsening restricted to pieces at or beyond the frontier. Each pass repairs
every mixed block and chooses, per workflow, the merge that adds the least
footprint-duration area, a proxy for how much additional lock holding the
merge causes. Each pass strictly reduces the number of pieces, so the loop
terminates, and by Theorem~\ref{thm:char} it must succeed once the sealing has
certified.

The same algorithm with an empty frontier is the launch-time certifier: it
takes the finest chopping the agents declare and returns the finest plan it
can certify.
